\clearpage
\section*{Exercise 8: Increasing the Amber LED Duration}
\addcontentsline{toc}{section}{Exercise 8: Increasing the Amber LED Duration}

\begin{questionbox}{1. Problem Statement}
Use Button 2 to modify the amber-light duration and Button 3 to confirm the selected value. The red and green timings must remain consistent with the selected amber duration.
\end{questionbox}

\vspace{0.4cm}
\textbf{2. Amber-duration Configuration}

Mode 3 loads \texttt{yellowTimeSeconds} into \texttt{selectedTimeSeconds}. Although the PDF calls this color amber, the source code uses the name \texttt{yellow}. A short Button 2 press increases the selected value by one unit; a long press requests a decrement. Button 3 saves the selected amber duration.

The confirmation logic updates the amber and red values while leaving the green duration unchanged:

\begin{lstlisting}[language=C, caption={Amber timing confirmation}, label={lst:exercise8_amber_confirm}, captionpos=b]
else if (mode == MODE_SET_YELLOW) {
    yellowTimeSeconds = selectedTimeSeconds;
    redTimeSeconds =
        (uint8_t)(greenTimeSeconds + yellowTimeSeconds);
}
\end{lstlisting}

The timing relationship remains
\[
    T_{\mathrm{red}} = T_{\mathrm{green}} + T_{\mathrm{amber}}.
\]
The upper bound is calculated as \(99 - T_{\mathrm{green}}\), so the derived red value does not exceed 99. The lower bound for the selected amber value is 1. At the current defaults, green is 7 and amber is 3; changing amber to 4 and confirming produces red = 11, amber = 4, and green = 7.

\textbf{3. Display and LED Feedback}

While Mode 3 is active, the vertical two-digit display shows the proposed amber duration and the horizontal display shows mode number 3. The amber/yellow logical outputs blink while red and green outputs remain off. The value only becomes the saved configuration when Button 3 is pressed.

\begin{resultbox}
Changing the amber duration leaves green unchanged and recalculates red as green plus amber. The software uses a variable named \texttt{yellowTimeSeconds} for the amber value.
\end{resultbox}

\vspace{0.2cm}
\vspace{0.15cm}
\noindent\textbf{GitHub source repository:} \href{https://github.com/haunguyen-2006/STM32LAB1}{STM32LAB1}.

% IMAGE SLOT: Add the screenshot described below and save it as ex08_amber_duration_mode3.png in the images/ folder.
% If the file is absent, a labeled placeholder is displayed; the document remains compilable.
\begin{figure}[H]
    \centering
    \IfFileExists{images/ex08_amber_duration_mode3.png}{%
        \includegraphics[width=0.98\textwidth]{images/ex08_amber_duration_mode3.png}%
    }{%
        \fbox{\parbox[c][4.4cm][c]{0.92\textwidth}{%
            \centering\textbf{IMAGE SLOT -- Mode 3 amber-duration configuration}\\[6pt]
            Insert a Proteus screenshot of Mode 3 showing the amber/yellow LEDs blinking, the selected amber timing value, and mode number 3.\\[6pt]
            \texttt{images/ex08\_amber\_duration\_mode3.png}%
        }}%
    }
    \caption{Mode 3: amber-duration configuration and LED feedback}
    \label{fig:exercise8_amber_mode_simulation}
\end{figure}


